\section{Diagram 2}\label{sec:d2}

Diagram~2 of \cite[Fig.~11]{AK} consists of the $z$-segment $12$ and the $w$-triangle $345$,
and every vertex carries a circle of the colour of its edges \cite[\S8.2, p.~572]{AK}.
Exchanging $z$ and $w$ and renumbering the bodies are symmetries of \eqref{eq:system}, so we
assume that a Diagram~2 germ has exactly this diagram.
Throughout this section $A,A'$ denote the pair bodies, $\{A,A'\}=\{1,2\}$, and $B,B'$ bodies of
the triangle $\{3,4,5\}$; $M=m_1+m_2$, $M_T=m_3+m_4+m_5$ and $m'=(m_3,m_4,m_5)$.  The letter $t$
is a chart coordinate; the parameter of a germ (called $t$ in Lemma~\ref{lem:germ}) is called
$s$, and $\ord$ is the order in $s$.  The letters $x,y,u$ are those of \cite[\S8.2]{AK}.

The proof has four steps.  We blow up \eqref{eq:system} along the diagram; every Diagram~2 germ
converges to a point of the boundary fibre $t=0$, a \emph{label}
(Proposition~\ref{prop:d2-cover}).  A torque identity of the pair forces $\Phi=0$ at the limit
label, which puts the masses on the loci $L_j$, $K_j$, $P_j$ of \cite[\S8.2]{AK}
(Proposition~\ref{prop:d2-first}).  Where the selected twenty equations define a graph over
$t$, the two remaining equations reduce to $tq(t)=0$, with $q(0)=\Phi$ and $q'(0)$ given by a
quadratic form of the triangle (Lemmas~\ref{lem:d2-isolation} and~\ref{lem:d2-second}).  At
fold labels without such a graph we instead use one- or two-parameter local reductions.
Finally every label with $\Phi=0$ at positive masses is shown not to be a limit of germs
(\S\ref{sec:d2-cases}), in most regular cases because $q'(0)\ne0$.  No reality condition is used.

\subsection{The chart and its labels}\label{sec:d2-chart}

The chart has coordinates $(t,Z,W,D)\in\C^{21}$, $Z=(Z_1,\dots,Z_5)$, $W=(W_1,\dots,W_5)$,
$D=(D_{kl})_{k<l}$, related to the coordinates of \eqref{eq:system} by
\begin{equation}\label{eq:d2-chart}
\begin{gathered}
  z_A=t^{-1}Z_A,\quad w_A=tW_A,\quad z_B=tZ_B,\quad w_B=t^{-1}W_B,\\
  \delta_{12}=D_{12},\quad \delta_{BB'}=D_{BB'},\quad \delta_{AB}=tD_{AB}.
\end{gathered}
\end{equation}
The single-index capitals $Z_k,W_k$ are chart coordinates, unrelated to the two-index quantities
$Z_{kl},W_{kl}$ used for diagrams.  Multiplying the coordinate equations of \eqref{eq:system} for
$z_A,w_A,z_B,w_B$ by $t,t^{-1},t^{-1},t$ gives the polynomials
\begin{equation}\label{eq:d2-rows}
\begin{aligned}
  F^Z_A&=Z_A-m_{A'}D_{12}^3(Z_A-Z_{A'})-t^3\textstyle\sum_Bm_BD_{AB}^3(Z_A-t^2Z_B),\\
  F^W_A&=W_A-m_{A'}D_{12}^3(W_A-W_{A'})+t\textstyle\sum_Bm_BD_{AB}^3(W_B-t^2W_A),\\
  F^Z_B&=Z_B-\textstyle\sum_{B'\ne B}m_{B'}D_{BB'}^3(Z_B-Z_{B'})
         +t\textstyle\sum_Am_AD_{AB}^3(Z_A-t^2Z_B),\\
  F^W_B&=W_B-\textstyle\sum_{B'\ne B}m_{B'}D_{BB'}^3(W_B-W_{B'})
         -t^3\textstyle\sum_Am_AD_{AB}^3(W_B-t^2W_A),
\end{aligned}
\end{equation}
and the distance equations become $T_{12}=D_{12}^2(Z_2-Z_1)(W_2-W_1)-1$,
$T_{BB'}=D_{BB'}^2(Z_{B'}-Z_B)(W_{B'}-W_B)-1$ ($B<B'$) and
$T_{AB}=D_{AB}^2(t^2Z_B-Z_A)(W_B-t^2W_A)-1$.  The gauge of \eqref{eq:system} is replaced by
$G_1=Z_2-Z_1-1$ and $G_2=W_5-1$.
Let $\mathcal X_m\subset\C^{21}$ be the zero set of these $22$ polynomials and
$\mathcal B_m=\mathcal X_m\cap\{t=0\}$; the points of $\mathcal B_m$ are the \emph{labels}.

\begin{definition}[Shape factors]\label{def:d2-shapes}
For $\varsigma\in\{\pm1\}^3$ let $N_\varsigma(w_3,w_4,w_5)$ be the numerator of
$w_3/(\varsigma_3(w_5-w_4)^2)+w_4/(\varsigma_4(w_3-w_5)^2)+w_5/(\varsigma_5(w_4-w_3)^2)$.  The
quintics $S_1,S_2,S_3$ are $N_\varsigma$ for $\varsigma=(1,-1,-1),(1,-1,1),(1,1,-1)$, and
$N_{(1,1,1)}$ is a constant multiple of $S_4S_5$, where
\begin{gather*}
  S_4=w_3^3+w_4^3+w_5^3-w_3^2w_4-w_3w_4^2-w_3^2w_5-w_3w_5^2-w_4^2w_5-w_4w_5^2+3w_3w_4w_5,\\
  S_5=w_3^2+w_4^2+w_5^2-w_3w_4-w_3w_5-w_4w_5 .
\end{gather*}
The quartics $S_6,S_7,S_8$ are the factors of
$\operatorname{Res}_p\bigl((w_4-w_3)^3-p^3(w_5-w_4)^3,\ (w_4-w_3)^3-(-1-p)^3(w_3-w_5)^3\bigr)$
(which is $S_5^3S_6S_7S_8$ up to a constant) whose coefficients of $w_3^3w_4$ are $-2$, $1$ and
$-5$.  These are the factors of the complex three-body central configuration locus of
\cite[Appendix, pp.~585--586]{AK}.
\end{definition}

\begin{lemma}[computer-assisted]\label{lem:d2-fibre}
Let $m\in\Rp^5$.  A point $p=(0,Z,W,D)$ is a label if and only if
\begin{alphalist}
  \item $Z_1=-m_2/M$, $Z_2=m_1/M$, $D_{12}^3=1/M$, $W_1=-m_2D_{12}$ and $W_2=m_1D_{12}$;
  \item $F^Z_B=F^W_B=T_{BB'}=0$ at $t=0$ (system \eqref{eq:system} for the bodies $3,4,5$,
  without its gauge) and $W_5=1$;
  \item $D_{AB}^2Z_AW_B=-1$ for all $A,B$.
\end{alphalist}
At a label $\sum_Bm_BZ_B=\sum_Bm_BW_B=0$, every $W_B\ne0$, the $W_B$ are pairwise distinct, and
the \emph{triangle} $(Z_B,W_B,D_{BB'})$ is of one of two kinds.
\begin{romanlist}
  \item \emph{Lagrange:} $D_{BB'}=M_T^{-1/3}\omega^{a_{BB'}}$, $\omega=e^{2\pi i/3}$,
  $a_{BB'}\in\{0,1,2\}$ ($27$ classes), and $\kappa=(W_4-W_3)/(W_5-W_4)$ is a root of
  $\omega^{a_{45}}\kappa^2+(\omega^{a_{34}}+\omega^{a_{45}}-\omega^{a_{35}})\kappa+\omega^{a_{34}}$.
  The root is double exactly for the six classes with $a_{34},a_{35},a_{45}$ distinct; these are
  the \emph{fold} triangles.  $(W_3,W_4,W_5)$ is a zero of $S_5$ if $a_{34}=a_{35}=a_{45}$ or the
  $a_{BB'}$ are distinct, of $S_6$ if $a_{34}=a_{35}\ne a_{45}$, of $S_7$ if
  $a_{34}=a_{45}\ne a_{35}$, and of $S_8$ if $a_{35}=a_{45}\ne a_{34}$.
  \item \emph{Euler:} $Z_B=c^{-2}W_B$ for some $c\in\C^*$, $D_{BB'}=e_{BB'}c/(W_{B'}-W_B)$ for
  $B<B'$ with $e_{BB'}=\pm1$, and $(W_3,W_4,W_5)$ is a zero of $S_1$, $S_2$, $S_3$ or $S_4S_5$
  according to the class of $(e_{34},e_{35},e_{45})$ modulo a global sign.  An Euler triangle on
  $S_5$ is a Lagrange fold triangle.
\end{romanlist}
\end{lemma}

\begin{proof}
At $t=0$ the crossed terms of \eqref{eq:d2-rows} vanish and $T_{AB}=-D_{AB}^2Z_AW_B-1$; this
gives (b), (c) and, with $G_1$, the system $Z_A=m_{A'}D_{12}^3(Z_A-Z_{A'})$,
$W_A=m_{A'}D_{12}^3(W_A-W_{A'})$, $T_{12}=0$, whose solutions are (a).  Summing $m_BF^Z_B$ and
$m_BF^W_B$ gives the centre-of-mass relations.  If $W_B=0$ then $T_{1B}=-1$, and if
$W_B=W_{B'}$ then $T_{BB'}=-1$.  If the vectors $(1)_B,(Z_B)_B,(W_B)_B$ are linearly independent,
then $D_{BB'}^3=1/M_T$ for all three edges.  Indeed, the linear map
$X\mapsto(\sum_{B'\ne B}m_{B'}D_{BB'}^3(X_B-X_{B'}))_B$ of $\C^3$ fixes $(Z_B)$ and $(W_B)$ by (b)
and kills $(1)_B$, and so does $X\mapsto X-M_T^{-1}(\sum_Bm_BX_B)(1)_B$, by the centre-of-mass
relations; as the three vectors form a basis, the two maps agree, and their entries at
$(B,B')$, $B\ne B'$, give $m_{B'}D_{BB'}^3=m_{B'}/M_T$.  Then the coordinate equations
reduce to the centre-of-mass relations, the ties give
$Z_{B'}-Z_B=M_T^{2/3}\omega^{a_{BB'}}/(W_{B'}-W_B)$, and $Z_5-Z_3=(Z_4-Z_3)+(Z_5-Z_4)$ is the
quadratic for $\kappa$.  Otherwise $Z=\alpha W+\beta$; the centre of mass gives $\beta=0$, the
ties give $\alpha\ne0$, and with $\alpha=c^{-2}$ they give the form of $D_{BB'}$.  For the double
roots write $c_{BB'}=\omega^{a_{BB'}}$: if two of the $c_{BB'}$ agree, the discriminant is
$b(b-4a)\ne0$ with $|a|=|b|=1$ ($a$ the repeated value, $b$ the remaining one), while if they are
distinct, $c_{34}+c_{45}=-c_{35}$ and $c_{34}c_{45}=c_{35}^2$, so the discriminant vanishes.  The
assignment of the shape factors is computer-assisted.
\end{proof}

The assignment is an exact computation over $\mathbb Q(\omega)$.  For each of the $27$ Lagrange
classes we reduce the factor $S_j$ of the class, as a polynomial in $\kappa$, modulo the quadratic
for $\kappa$, and for a fold class we evaluate it at the double root; for each of the eight sign
vectors $(e_{34},e_{35},e_{45})$ we identify the corresponding $N_\varsigma$ with one of $\pm S_1$,
$\pm S_2$, $\pm S_3$, $\pm S_4S_5$; and for the Euler triangles on $S_5$ we check that the three
numbers $D_{BB'}^3$ coincide modulo $S_5$ (Appendix~\ref{app:verification}).

\begin{definition}[Mass polynomials]\label{def:d2-mass}
For positions $W_3,W_4,W_5$ put $f_B=m_B^2/W_B$,
$\Lambda=f_3^2+f_4^2+f_5^2-2f_3f_4-2f_3f_5-2f_4f_5$ and
$\Theta=3(f_3^2+f_4^2+f_5^2)+2f_3f_5+2f_4f_5-6f_3f_4$.  Equations (25) and (28) of
\cite[\S8.2]{AK} are
\begin{equation}\label{eq:d2-25-28}
  \Lambda=0,\qquad
  \Lambda\,y^2+2(f_3-f_4+f_5)(f_3-f_4-f_5)\,y+\tfrac12(\Theta-\Lambda)=0 .
\end{equation}
For a rational expression $H$ in $W_3,W_4,W_5$ whose denominator is a monomial, let $H^\flat$ be
the polynomial in $W_4$ obtained by clearing the position denominators, substituting $W_5=1$ and
$W_3=-(m_4W_4+m_5)/m_3$ (centre of mass), and clearing the resulting power of $m_3$.  Define
$L_j(m')$ and $P_j(y)$ by
\[
  \operatorname{Res}_{W_4}(S_j^\flat,\Lambda^\flat)=c_j\,L_j,\qquad
  \operatorname{Res}_{W_4}\bigl(S_j^\flat,(28)^\flat\bigr)=m_3^{2\deg S_j}P_j(y),
\]
where $c_j$ has no zero on $\Rp^3$; for $j=5$ the first resultant is
$m_3^4M_T^2L_5L_5'$.  Write $P_j(y)=E_j(y^2)+yO_j(y^2)$ and $K_j=P_j(0)=E_j(0)$.
\end{definition}

These are the factors $L_1,\dots,L_8,L_5'$ and the polynomials $P_1,\dots,P_8$ of
\cite[pp.~573--575]{AK}, reconstructed from the appendix factors; the displayed coefficients and
published symmetries of \cite{AK} are reproduced.  In particular
$2L_5=(m_3-m_4)^2+(m_4-m_5)^2+(m_5-m_3)^2$, and $L_5'$ has degree $12$ and $91$ coefficients, all
positive.
Permutations of $(m_3,m_4,m_5)$ permute $\{L_1,L_2,L_3\}$ and $\{L_6,L_7,L_8\}$ up to nonzero
constants; the exchange of bodies $3$ and $4$ sends $S_1$ to $-S_2$, exchanges $P_1\leftrightarrow
P_2$ and $P_6\leftrightarrow P_7$, and fixes $P_3,P_4,P_5,P_8$.
These statements are proved by exact computation over $\mathbb Q$, with the masses and $y$ as
indeterminates: the resultants defining $L_j$ and $L_5'$ are computed and factored, those
defining $P_j$ are determined by evaluation on unisolvent grids of points modulo three primes,
using a priori bounds on their degrees and coefficients, and the decompositions $P_j=E_j+yO_j$,
the leading coefficients and the symmetries are checked by expansion and substitution
(Appendix~\ref{app:verification}).
\subsection{Coverage}\label{sec:d2-cover}

\begin{proposition}[Coverage]\label{prop:d2-cover}
Let $m\in\Rp^5$ and let $\gamma$ be a Diagram~2 germ on a curve in $A_m$.  Along $\gamma$ put
$t=(z_{12}w_5)^{-1/2}$ \textup{(}either branch\textup{)} and $\lambda=(tz_{12})^{-1}$, and let
$(Z,W,D)$ be given by \eqref{eq:d2-chart} applied to the rotated configuration
$(\lambda z,\lambda^{-1}w,\delta)$.  Then $t,Z,W,D$ are Puiseux series in $s$ with $\ord t>0$ and
nonnegative orders, $(t,Z,W,D)\in\mathcal X_m$ for $s\ne0$, and the limit $p$ as $s\to0$ is a label at
which $Z_1,Z_2$, $W_3,W_4,W_5$ and all $D_{kl}$ are nonzero.
\end{proposition}

\begin{proof}
The equations of \eqref{eq:system} other than its gauge are invariant under the rotations
$(z,w)\mapsto(\lambda z,\lambda^{-1}w)$, the chart rows are these equations times powers of $t$,
and $G_1=\lambda tz_{12}-1=0$, $G_2=\lambda^{-1}tw_5-1=t^2z_{12}w_5-1=0$; so the chart point lies
in $\mathcal X_m$.  Rotating $\gamma$ beforehand changes neither $t$ nor the chart point, so for the
orders we may normalize as in Lemma~\ref{lem:germ} and \cite[p.~547]{AK}: the maximal components
of $\mathcal Z$ and $\mathcal W$ have one order, and we use units in which it is $-2$.  By
Fact~\ref{fact:rules}:
\begin{enumerate}
  \item The pair is an isolated component of the $z$-diagram with $z$-circles, so its $z$-stroke
  is maximal (Rule~1e): $\ord z_{12}=-2$; by Estimate~2, $r_{12}\approx1$, so $\ord w_{12}=2$ and
  $\ord\delta_{12}=0$.
  \item The triangle is an isolated $w$-component with $w$-circles; it has a maximal $w$-stroke
  (Rule~1e), and by Rule~2e all three are maximal: $\ord w_{BB'}=-2$, $\ord z_{BB'}=2$,
  $\ord\delta_{BB'}=0$.
  \item The circles give $\ord z_A=\ord w_B=-2$.  The bodies $3,4,5$ carry no $z$-stroke and
  $1,2$ no $w$-stroke, hence no circle of that colour (Rule~1a): $\ord z_B>-2$, $\ord w_A>-2$.
  \item Hence $\ord z_{AB}=\ord w_{AB}=-2$ and $\ord\delta_{AB}=2$.
  \item In $w_A=\sum_{l\ne A}m_l\delta_{Al}^3(w_A-w_l)$ the pair term has order $2$ and the
  crossed terms order $4$, so $\ord w_A\ge2$; likewise $\ord z_B\ge2$.
  \item So $\ord t=2$ and $\ord\lambda=0$; $Z_A=\lambda tz_A$, $W_B=\lambda^{-1}tw_B$, $D_{12}$,
  $D_{BB'}$ and $D_{AB}=\delta_{AB}/t$ have order exactly $0$, and $W_A=\lambda^{-1}w_A/t$,
  $Z_B=\lambda z_B/t$ have order $\ge0$.
\end{enumerate}
Hence every chart coordinate converges, $t\to0$, and the coordinates of order exactly $0$ have
nonzero limits.  The limit lies in the closed set $\mathcal X_m$ and has $t=0$.
\end{proof}

\subsection{First conditions}\label{sec:d2-first}

\begin{lemma}[Syzygies]\label{lem:d2-syzygy}
Identically in $(t,Z,W,D)$,
\begin{gather*}
  \textstyle\sum_{k=1}^5m_k(W_kF^Z_k-Z_kF^W_k)=0,\qquad
  \sum_Am_A(W_AF^Z_A-Z_AF^W_A)=t\,Q,\\
  Q=\textstyle\sum_{A,B}m_Am_BD_{AB}^3(t^4W_AZ_B-Z_AW_B).
\end{gather*}
At $t=0$ the only rows with a nonzero $t$-derivative are $F^W_A$ and $F^Z_B$, with
$\partial_tF^W_A=\sum_Bm_BD_{AB}^3W_B$ and $\partial_tF^Z_B=\sum_Am_AD_{AB}^3Z_A$.  At a label,
$Q=\Phi:=\sum_Bm_B\varphi_B$ with $\varphi_B=m_1D_{1B}+m_2D_{2B}$, and
$\partial_tF^Z_B=-\varphi_B/W_B$.
\end{lemma}

\begin{proof}
Since $W_AF^Z_A=w_AF^z_A$, $Z_AF^W_A=z_AF^w_A$, $W_BF^Z_B=w_BF^z_B$ and $Z_BF^W_B=z_BF^w_B$ for
the coordinate equations $F^z_k,F^w_k$ of \eqref{eq:system}, the first identity is the torque
identity $\sum_km_k(w_kF^z_k-z_kF^w_k)=0$, which holds by antisymmetry.  In the second sum the
terms with $D_{12}$ cancel by antisymmetry and the crossed terms give $tQ$.  The derivatives are
read off from \eqref{eq:d2-rows}, and Lemma~\ref{lem:d2-fibre}(c) gives
$D_{AB}^3Z_AW_B=-D_{AB}$ and $D_{AB}^3Z_A=-D_{AB}/W_B$.
\end{proof}

\emph{Crossed data.}  At a label fix $v_B$ with $v_B^2=1/W_B$.  By Lemma~\ref{lem:d2-fibre}(a),(c),
$D_{1B}^2=M/(m_2W_B)$ and $D_{2B}^2=-M/(m_1W_B)$, so $D_{1B}=s_B\sqrt{M/m_2}\,v_B$ and
$D_{2B}=is'_B\sqrt{M/m_1}\,v_B$ with $s_B,s'_B\in\{\pm1\}$ (positive square roots); these are the
$64$ crossed-sign assignments.  Put $\sigma_B=s_Bs'_B$,
$c_\pm=\sqrt M\,(m_1/\sqrt{m_2}\pm i\,m_2/\sqrt{m_1})$ and $x=i(m_2/m_1)^{3/2}$.  Then
$c_+c_-=M(m_1^3+m_2^3)/(m_1m_2)>0$, $\varphi_B=c_{\sigma_B}s_Bv_B$, and
$m_2D_{2B}/(m_1D_{1B})=\sigma_Bx$ with $x^2=-m_2^3/m_1^3$.  The label is \emph{same-sign} if
$\sigma_3=\sigma_4=\sigma_5$ and \emph{one-different} otherwise (two and six of the eight
patterns); then the \emph{odd body} is the $B$ whose $\sigma_B$ differs from the majority sign
$\sigma$.  In both cases put $\tilde v_B=\varphi_B/c_\sigma$, so that $\tilde v_B=s_Bv_B$ if
$\sigma_B=\sigma$ and $\tilde v_k=\rho\,s_kv_k$ for the odd body $k$, where
$\rho=c_{-\sigma}/c_\sigma$ satisfies $|\rho|=1$, $\rho\ne\pm1$.  Thus
$\Phi=c_\sigma\sum_Bm_B\tilde v_B$.
\begin{proposition}[First conditions]\label{prop:d2-first}
Let $m\in\Rp^5$ and let $p$ be the limit label of a Diagram~2 germ
\textup{(}Proposition~\ref{prop:d2-cover}\textup{)}.  Then $\Phi(p)=0$, and:
\begin{alphalist}
  \item if $p$ is same-sign, its triangle satisfies \textup{(25)}, $\Lambda=0$, and
  $L_j(m')=0$ for every $j\ne5$ with $S_j(W_3,W_4,W_5)=0$ \textup{(}$L_5L_5'=0$ if $j=5$\textup{)};
  \item if $p$ is one-different with odd body $5$ and majority sign $\sigma$, its triangle satisfies
  \textup{(28)} with $y=\sigma(x+x^{-1})/2$, and $u:=y^2=-(T-1)^2/(4T)$, $T=(m_1/m_2)^3$.  Hence
  $P_j(y)=0$ for every $j$ with $S_j(W_3,W_4,W_5)=0$: if $m_1=m_2$ then $K_j(m')=0$, and if
  $m_1\ne m_2$ then $u<0$ and $E_j(u)=O_j(u)=0$.
\end{alphalist}
\end{proposition}

\begin{proof}
Along the germ all rows vanish and $t\not\equiv0$, so Lemma~\ref{lem:d2-syzygy} gives $Q\equiv0$
along it, and $\Phi(p)=Q(p)=0$ by continuity.  No Jacobian rank is used, so this holds at every
label, including fold labels.
In case (a), $\Phi=c_\sigma\sum_Bm_Bs_Bv_B$ with $c_\sigma\ne0$, and the product of
$m_3v_3\pm m_4v_4\pm m_5v_5$ over the four sign choices is $\Lambda$, since
$(m_Bv_B)^2=f_B$.  In case (b) put $r_B=m_1m_BD_{1B}$, so $r_B^2=(m_1^2M/m_2)f_B$, and
$x'=\sigma x$; then $\Phi=(r_3+r_4)(1+x')+r_5(1-x')$, which is (26) of \cite[\S8.2]{AK}, and the
product over the four sign changes of $r_3,r_4$ is $4x'^2(m_1^2M/m_2)^2$ times the left side of
(28) with $y=(x'+1/x')/2$.
In both cases $W_4$ is a common root of $S_j^\flat$ and $\Lambda^\flat$, resp.\ $(28)^\flat$, so
the resultants of Definition~\ref{def:d2-mass} vanish, and $c_j\ne0$, $m_3\ne0$.  Finally
$y=iv$ with $v=\sigma\bigl((m_2/m_1)^{3/2}-(m_1/m_2)^{3/2}\bigr)/2\in\R$, which vanishes iff
$m_1=m_2$, and $E_j,O_j$ have real coefficients; so $P_j(iv)=E_j(-v^2)+ivO_j(-v^2)=0$ gives the
last assertion.
\end{proof}

\subsection{The isolation mechanism}\label{sec:d2-isolation}

Let $F_{\mathrm{sel}}$ denote the twenty rows other than $F^Z_2$ and $F^Z_5$.  A label $p$ is
\emph{excluded} if no holomorphic curve germ $c\colon(\C,0)\to(\mathcal X_m,p)$ has
$t\circ c\not\equiv0$; an isolated point of $\mathcal X_m$ is excluded.

\begin{lemma}[Local analytic lemma]\label{lem:d2-isolation}
Let $m\in\Rp^5$ and $p\in\mathcal B_m$.
\begin{alphalist}
  \item If $p$ is excluded, it is not the limit label of a Diagram~2 germ.
  \item Suppose $dF_{\mathrm{sel}}(p)$ has rank $20$, so that near $p$ the set
  $\{F_{\mathrm{sel}}=0\}$ is a smooth curve $C$.  Then near $p$ the set $\mathcal X_m$ consists of the
  points of $C$ at which $tQ=0$, and $p$ is excluded if $Q|_C\not\equiv0$.  If moreover
  $\partial F_{\mathrm{sel}}/\partial(Z,W,D)(p)$ is invertible, then $C$ is the graph of an
  analytic map $t\mapsto\phi(t)$, and $p$ is an isolated point of $\mathcal X_m$ as soon as
  $q(t)=Q(t,\phi(t))\not\equiv0$, in particular if $q(0)=\Phi(p)\ne0$ or $q'(0)\ne0$.
  \item Suppose some twenty of the twenty-two polynomials have a differential of rank $20$ at
  $p$, so that near $p$ their common zeros form a smooth curve $C'$.  If one of the other two
  polynomials does not vanish identically on $C'$, then $p$ is an isolated point of $\mathcal X_m$.
\end{alphalist}
\end{lemma}

\begin{proof}
(a) By Proposition~\ref{prop:d2-cover}, a Diagram~2 germ with limit $p$ gives, after the
substitution $s=\tau^N$ that clears the denominators of the Puiseux exponents, a holomorphic germ
$\tau\mapsto(t,Z,W,D)$ in $\mathcal X_m$ at $p$ with $t\not\equiv0$.
(b) On $C$ all rows except $F^Z_2,F^Z_5$ vanish and $W_5=1$, so Lemma~\ref{lem:d2-syzygy} gives
$m_2W_2F^Z_2=tQ=-m_5F^Z_5$ on $C$; and $W_2(p)=m_1D_{12}(p)\ne0$.  This describes $\mathcal X_m$
near $p$.  A germ $c$ in $\mathcal X_m$ at $p$ lies in $C$, and if $t\circ c\not\equiv0$ then
$Q\circ c\equiv0$; when
$Q|_C\not\equiv0$ its zeros on the curve $C$ are isolated, so $c$ is constant, which contradicts
$t\circ c\not\equiv0$.  In the graph case $\mathcal X_m$ is, near $p$, the set of points $(t,\phi(t))$ with
$tq(t)=0$, which is $\{p\}$ when $q\not\equiv0$.
(c) Near $p$, $\mathcal X_m$ lies in $C'$ and in the zero set of the other polynomial, whose restriction
to the smooth curve $C'$ is analytic and not identically zero, hence has an isolated zero at $p$.
\end{proof}

At a few labels the Jacobian of the twenty-two polynomials has rank $19$.  There we use the
following formal obstruction, a Newton--Puiseux argument on the blow-up of a Lyapunov--Schmidt
reduction.

\begin{lemma}[Corank-two isolation]\label{lem:d2-corank-two}
Let $m\in\Rp^5$ and $p\in\mathcal B_m$.  Let $F_R$ be nineteen of the twenty-two chart polynomials,
$F_{R'}$ the other three, and $y\ne t$ a chart coordinate such that the derivative of
$(F_R,t,y)\colon\C^{21}\to\C^{21}$ at $p$ is invertible.  By the implicit function theorem,
near $p$ the set $\{F_R=0\}$ is the image of an injective analytic map $\phi$, defined near
$0\in\C^2$, with $\phi(0)=p$ and $(y,t)\circ\phi(a,t)=(y(p)+a,t)$, and near $p$
\[
  \mathcal X_m=\phi\bigl(g^{-1}(0)\bigr),\qquad g=F_{R'}\circ\phi\colon(\C^2,0)\to(\C^3,0).
\]
Let $g^{(j)}$ be the homogeneous part of degree $j$ of $g$, and for $\xi\in\C^2$ let
$\Pi_\xi=Dg^{(2)}(\xi)$, a $3\times2$ matrix.  Suppose:
\begin{alphalist}
  \item $g^{(1)}=0$; equivalently, the Jacobian at $p$ of all twenty-two polynomials has
  rank~$19$;
  \item the components of $g^{(2)}$ have only finitely many common zeros $[\xi_1],\dots,[\xi_r]$
  in $\mathbb P^1$, and $\Pi_{\xi_j}\ne0$ for every $j$;
  \item for every $j$ there is $N\ge3$ such that the vector $r_N(\xi_j)$ defined below exists
  and does not lie in $\Pi_{\xi_j}\C^2$.
\end{alphalist}
Then no nonconstant formal arc $x\in\C[[\tau]]^2$ with $x(0)=0$ satisfies $g(x)=0$.  Hence $p$
is an isolated point of $\mathcal X_m$; in particular $p$ is excluded.  The same holds if
\textup{(a)--(c)} hold for $Kg(L\,\cdot)$, $K\in GL_3(\C)$, $L\in GL_2(\C)$, in place
of $g$.

The vectors $r_N(\xi)$ are defined for a common zero $\xi$ of $g^{(2)}$ and a fixed
$e\in\C^2\setminus\C\xi$.  As $\Pi_\xi\xi=2g^{(2)}(\xi)=0$ and $\Pi_\xi\ne0$, the vector
$\varpi=\Pi_\xi e$ spans the line $\Pi_\xi\C^2$.  Put $\beta_1=0$.  For $n\ge1$, if $\beta_n$ is
defined, let $r_{n+2}(\xi)$ be the coefficient of $\lambda^{n+2}$ in
$g(\lambda\xi+\beta_n(\lambda)e)$; if moreover $r_{n+2}(\xi)=-b_{n+1}\varpi$ with $b_{n+1}\in\C$,
put $\beta_{n+1}(\lambda)=\beta_n(\lambda)+b_{n+1}\lambda^{n+1}$.  Thus
\[
  r_3(\xi)=g^{(3)}(\xi),\qquad
  r_4(\xi)=g^{(4)}(\xi)+b_2\,Dg^{(3)}(\xi)e+b_2^2\,g^{(2)}(e).
\]
Condition (c) with $N=3$ says that the $3\times3$ matrix $\bigl(\Pi_{\xi_j}\,|\,r_3(\xi_j)\bigr)$
has rank two; with $N=4$ it says that $\bigl(\Pi_{\xi_j}\,|\,r_3(\xi_j)\bigr)$ has rank one and
$\bigl(\Pi_{\xi_j}\,|\,r_4(\xi_j)\bigr)$ has rank two.
\end{lemma}

\begin{proof}
The points of $\{F_R=0\}$ near $p$ are the points $\phi(a,t)$, and $\phi(a,t)\in\mathcal X_m$ if
and only if $g(a,t)=0$.  If $p$ were not isolated in $\mathcal X_m$, the curve selection lemma
\cite[Thm.~2.5.5]{BCR}, applied to the real algebraic set $\mathcal X_m\subset\C^{21}=\R^{42}$,
would give a nonconstant semialgebraic arc $\gamma\colon[0,\varepsilon)\to\mathcal X_m$ with
$\gamma(0)=p$.  Its
coordinates are convergent Puiseux series, so $x=(y\circ\gamma-y(p),t\circ\gamma)$ is, after a
substitution $s=\tau^{N'}$, a nonconstant arc in $\C[[\tau]]^2$ with $g(x)=0$.  So the first
assertion implies the second.  The zero set of $Kg(L\,\cdot)$ is $L^{-1}g^{-1}(0)$, which gives
the last assertion.

Let $x$ be a nonconstant formal arc with $x(0)=0$ and $g(x)=0$, and write
$x=\tau^k\xi'+O(\tau^{k+1})$ with $\xi'\ne0$.  By (a),
$g(x)=\tau^{2k}g^{(2)}(\xi')+O(\tau^{2k+1})$, so $\xi'\in\C^*\xi$ for one of the $\xi=\xi_j$
of (b).  Write $x=u\xi+\tilde\theta e$; then $\ord u=k<\ord\tilde\theta$.  Take $\lambda=u(\tau)$ as
the new parameter; then $\tau$ is a power series in $\lambda^{1/k}$ with nonzero linear term,
and $x=\lambda\xi+\theta e$ with $\theta\in\C[[\lambda^{1/k}]]$ and $\ord\theta>1$ ($\theta=0$ is
allowed).

We show by induction on $n\ge1$ that $\beta_n$ is defined, that
$g(\lambda\xi+\beta_ne)=O(\lambda^{n+2})$ and that $o_n=\ord(\theta-\beta_n)>n$.  For $n=N-2$
this contradicts (c).  For $n=1$ the claim holds by (a), $g^{(2)}(\xi)=0$ and $\ord\theta>1$.
Taylor's formula in the direction $e$ gives
\[
  0=g(\lambda\xi+\theta e)
   =g(\lambda\xi+\beta_ne)+(\theta-\beta_n)\bigl(\lambda\varpi+O(\lambda^2)\bigr)
    +O(\lambda^{2o_n}),
\]
because $Dg(\lambda\xi+\beta_ne)e=\lambda\Pi_\xi e+\beta_n\Pi_ee+O(\lambda^2)
=\lambda\varpi+O(\lambda^2)$ by (a) and $\ord\beta_n\ge2$.  If $o_n<n+1$, the term
$(\theta-\beta_n)\lambda\varpi$ has order $o_n+1<n+2$, every other term has larger order
($2o_n>o_n+1$), and its leading coefficient is a nonzero multiple of $\varpi\ne0$; this is
impossible.  Hence $o_n\ge n+1$, and the coefficient of $\lambda^{n+2}$ gives
$r_{n+2}(\xi)+b\varpi=0$, where $b$ is the coefficient of $\lambda^{n+1}$ in $\theta$.  So
$\beta_{n+1}$ is defined, $b_{n+1}=b$ since $\varpi\ne0$, and $o_{n+1}>n+1$.  The same expansion
with $\beta_{n+1}$ in place of $\theta$ gives
$g(\lambda\xi+\beta_{n+1}e)=(r_{n+2}(\xi)+b_{n+1}\varpi)\lambda^{n+2}+O(\lambda^{n+3})
=O(\lambda^{n+3})$.
\end{proof}

\begin{remark}\label{rem:d2-corank-two}
The $\xi_j$ are the points where the strict transform of $g^{-1}(0)$ could meet the exceptional
curve of the blow-up of $0\in\C^2$, and the induction is the Newton--Puiseux algorithm in the
blow-up chart at $\xi_j$.  Arcs inside the boundary fibre $\{t=0\}$ need no separate argument:
they are tangent to the line $t=0$, which is then one of the $\xi_j$, and the same induction
excludes them.
\end{remark}

\emph{The triangle block and the form $Q_T$.}  For a label $p$ let $J_T(p)$ be the Jacobian at
$p$ of the nine rows $F^Z_3,F^W_3,F^Z_4,F^W_4,F^W_5,T_{34},T_{35},T_{45},G_2$ with respect to the
nine triangle coordinates $Z_B,W_B,D_{BB'}$.  If $J_T(p)$ is invertible and $\xi\in\C^3$, let
$(\zeta,\eta,\vartheta)$ solve $J_T(p)(\zeta,\eta,\vartheta)=r(\xi)$, where $r(\xi)$ has the
entries $\xi_3/W_3$ and $\xi_4/W_4$ in the rows $F^Z_3$, $F^Z_4$ and $0$ elsewhere, and put
\[
  Q_T(\xi)=-\tfrac12\textstyle\sum_Bm_B\,\xi_B\,\eta_B/W_B .
\]
$Q_T$ is a quadratic form determined by $m'$ and the triangle of $p$.

\begin{lemma}[Second-order coefficient]\label{lem:d2-second}
Let $p$ be a label with $\Phi(p)=0$ and $J_T(p)$ invertible.  Then
$\partial F_{\mathrm{sel}}/\partial(Z,W,D)(p)$ is invertible and, in the notation of
Lemma~\ref{lem:d2-isolation} and of the crossed data,
\[
  q(0)=0,\qquad q'(0)=c_\sigma^2\,Q_T(\tilde v).
\]
\end{lemma}

\begin{proof}
At $t=0$ the pair rows involve only $Z_A,W_A,D_{12}$, the triangle rows only the triangle
coordinates, and $T_{AB}$ involves no crossed coordinate other than $D_{AB}$, with
$\partial T_{AB}/\partial D_{AB}=-2D_{AB}Z_AW_B=2/D_{AB}$; so the Jacobian is block lower
triangular.  The pair block (rows $F^Z_1,F^W_1,F^W_2,T_{12},G_1$) is invertible for positive
masses (computer-assisted).  Differentiate $F_{\mathrm{sel}}(t,\phi(t))=0$ at $t=0$, using
Lemma~\ref{lem:d2-syzygy}.  The pair rows are forced by
$\partial_tF^W_A=-Z_A^{-1}\sum_Bm_BD_{AB}$, and
$\partial_tF^W_1-\partial_tF^W_2=M\Phi/(m_1m_2)=0$; for equal forcings the solution has
$Z_A'(0)=D_{12}'(0)=0$ (computer-assisted).  The triangle rows give
$J_T(p)(Z',W',D')_T(0)=r(\varphi)$, hence $W_B'(0)=c_\sigma\eta_B(\tilde v)$.  On $C$ the ties
$T_{AB}=0$ give $Q=\sum_{A,B}m_Am_BD_{AB}+O(t^2)$ and
$D_{AB}'=-\frac12D_{AB}(Z_A'/Z_A+W_B'/W_B)$.  Therefore
$q'(0)=\sum_{A,B}m_Am_BD'_{AB}(0)=-\frac12\sum_Bm_B\varphi_BW_B'(0)/W_B=c_\sigma^2Q_T(\tilde v)$.
\end{proof}

The two computer-assisted steps are exact computations with the masses as indeterminates, at
the pair part of Lemma~\ref{lem:d2-fibre}(a) and modulo $MD_{12}^3=1$: the determinant of the pair
block is a power of $D_{12}$ times factors of fixed sign for positive masses, and for equal
forcings the pair rows are solved with $Z_A'(0)=D_{12}'(0)=0$ (Appendix~\ref{app:verification}).

\begin{definition}\label{def:d2-psi}
For a label $p$ with $J_T(p)$ invertible put $\Psi_T(p)=Q_T(\tilde v)$ if $p$ is same-sign and
$\Psi_1(p)=Q_T(\tilde v)$ if $p$ is one-different.  The label is \emph{isolated at second order}
if $\Phi(p)=0$, $J_T(p)$ is invertible and $Q_T(\tilde v)\ne0$; by
Lemmas~\ref{lem:d2-isolation} and~\ref{lem:d2-second} it is then an isolated point of $\mathcal X_m$, and
the omitted rows $F^Z_2,F^Z_5$ vanish to order exactly two along the curve $C$.
\end{definition}

At a label with $\Phi(p)=0$ the vector $\tilde v$ is determined by the triangle and the signs
$s_B$ (for an odd body $k$, $m_k\tilde v_k=-\sum_{B\ne k}m_B\tilde v_B$), so $\Psi_T$ and
$\Psi_1$ do not involve $m_1,m_2$.  For a one-different label with odd body $5$ the condition
$\Phi=0$ requires $|m_3\tilde v_3+m_4\tilde v_4|=m_5|v_5|$, i.e.\ the vanishing of the
\emph{sheet function} $G=|m_3\tilde v_3+m_4\tilde v_4|^2-m_5^2$ (recall $W_5=1$), which again
depends on the triangle only; then $\tilde v_5=\rho s_5v_5$ fixes $\rho$ up to sign, hence
$m_1/m_2$ up to exchanging $m_1$ and $m_2$.

\subsection{The case analysis}\label{sec:d2-cases}

By Proposition~\ref{prop:d2-first} and Lemma~\ref{lem:d2-isolation}(a) it remains to show that
every label $p$ with $\Phi(p)=0$ at positive masses is excluded.  Table~\ref{tab:d2-cases} lists
the cases; the computer-assisted statements follow.  Ball arithmetic is IEEE arithmetic with
proved rounding bounds (Appendix~\ref{app:verification}).

\begin{table}[ht]
\centering
\footnotesize
\setlength{\tabcolsep}{3pt}
\begin{tabular}{p{2.5cm}p{2.5cm}p{3.3cm}p{3.4cm}}
\hline
Labels & Mass locus & Method & Computation \\
\hline
$\Phi\ne0$ (any) & all $m$ & first order, $q(0)=\Phi$ & exact identities \\
same-sign, $S_4$, or $S_5$ with $L_5'=0$ & $L_4=0$, $L_5'=0$ & no point in $\Rp^3$ & exact \\
same-sign, $S_5$ & $L_5=0$: $m_3=m_4=m_5$ & order three; corank-two blow-up & exact jets \\
same-sign, $S_1$--$S_3$, $S_6$--$S_8$ & rays of $L_1$--$L_3$, $L_6$--$L_8$ & second order,
$|\Psi_T|>0.3$ & Bernstein subdivision, Krawczyk method; interval arithmetic \\
\hline
one-different, $m_1=m_2$, $S_3,S_6$--$S_8$ & $K_3,K_6,K_7,K_8=0$ & no point in $\Rp^3$ &
Bernstein coefficients \\
one-different, $m_1=m_2$, $S_1,S_2$ & $K_1=0$, $K_2=0$ & second order & exact; interval cover \\
one-different, $m_1=m_2$, $S_4$ & $K_4=0$: four arcs & second order & exact \\
one-different, $m_1=m_2$, $S_5$ & $K_5=0$: $\Rp(1,1,\rho_*)$ & orders one to four &
exact; interval arithmetic \\
\hline
one-different, $m_1\ne m_2$, not fold & $E_j(u)=O_j(u)=0$ & second order, $\Psi_1\ne0$: ball
covers, exact lemmas & ball and interval arithmetic; exact \\
one-different, $m_1\ne m_2$, fold & $E_5(u)=O_5(u)=0$ & corank one; $Q|_C$ has a simple zero &
exact \\
\hline
\end{tabular}
\caption{The Diagram~2 cases.  A label with $\Phi=0$ lies on $S_j$ for the $j$ given in the
first column (Lemma~\ref{lem:d2-fibre}); one-different labels are taken with odd body~$5$.}
\label{tab:d2-cases}
\end{table}

\begin{lemma}[computer-assisted]\label{lem:d2-empty}
\textup{(a)} The polynomials $L_4$, $L_5'$, $K_3$, $K_6$, $K_7$ and $K_8$ have no zero in $\Rp^3$.
More precisely, at positive masses every $S_4$ triangle is real \textup{(}the centre-of-mass
section of $S_4$ is a cubic with positive discriminant\textup{)} and violates \textup{(25)}.
\textup{(b)} If $m_3=m_4=m_5$, no $S_j$ with $j\ne5$ has a root satisfying \textup{(25)}.
\end{lemma}
The proof is an exact computation over $\mathbb Q$.  For $L_4$ we use a rational parametrization
of $S_4$; on a real $S_4$ triangle at positive masses the positions $W_B$ do not all have the same
sign, and then $\Lambda$ is a sum of nonnegative terms, one of them positive.  $L_5'$ has
positive coefficients, and $K_3$, $K_6$, $K_7$, $K_8$ have nonnegative Bernstein coefficients on
twelve projective charts, three for each of them, each split into the two sectors of a blow-up
at a corner of the cone.  For (b) we check that $S_j^\flat$ and $\Lambda^\flat$ are coprime at
$m_3=m_4=m_5$ for $j\ne5$ (Appendix~\ref{app:verification}).

\begin{lemma}[computer-assisted]\label{lem:d2-L5}
Let $m\in\Rp^5$ with $m_3=m_4=m_5$.  Every label at $m$ whose triangle lies on $S_5$ \textup{(}an
all-equal Lagrange class or a fold triangle\textup{)} and which satisfies $\Phi=0$ or is a fold
label is an isolated point of $\mathcal X_m$.  At each such $m$ these are $72$ labels of the all-equal
classes and $1152$ fold labels, $1224$ in all.
\end{lemma}

\emph{Masses.}  Common mass scaling is an exact isomorphism of the gauged chart, so we may take
$m_3=m_4=m_5=1$; then $\delta_*:=M_T^{-1/3}=3^{-1/3}$.  Every positive pair $(m_1,m_2)$ is, for
exactly one $(\nu,M)\in(0,1)\times\Rp$,
\[
  m_1=\frac{4M\nu^2}{(1+\nu^2)^2},\qquad m_2=\frac{M(1-\nu^2)^2}{(1+\nu^2)^2},
\]
namely $M=m_1+m_2$ and $\nu=1/(\chi+\sqrt{1+\chi^2})$, where
$\chi=\sqrt{m_2/m_1}=(1-\nu^2)/(2\nu)$.

\emph{The all-equal labels.}  The $72$ labels of the all-equal classes have $\Phi=0$, an
invertible $J_T(p)$ and $q(0)=q'(0)=0$, and along $C$ we have $F^Z_5=\kappa_3t^3+O(t^4)$, so
$q''(0)=-2\kappa_3$.  For the nine labels of one sign class with $D_{BB'}=\omega^k\delta_*$,
$(W_3,W_4,W_5)=(\omega,\omega^2,1)$ and $D_{12}=\omega^jM^{-1/3}$,
\[
  \kappa_3=\omega^jM^{5/3}R_1(1-i\chi)+3\omega^k\delta_*MR_2(1-i\chi^5),\qquad
  R_1=\frac{18\nu^2(1-\nu^2)}{(1+\nu^2)^3},\quad R_2=\frac{6\nu^2(1+\nu^2)}{(1-\nu^2)^3}.
\]
These nine are the images of the case $j=k=0$ under the automorphisms
$\delta_*\mapsto\omega^k\delta_*$, $M^{1/3}\mapsto\omega^{-j}M^{1/3}$ of the coefficient field,
which fix the masses and the polynomials; the other $63$ labels are images of these nine under
the exchange of bodies $3$ and $4$, conjugation and the reversal of $t$.  Here $R_1,R_2>0$, and
$(1-i\chi)/(1-i\chi^5)$ is a positive multiple of $1+\chi^6+i(\chi^5-\chi)$, where
$|\chi^5-\chi|\le\chi+\chi^5\le1+\chi^6$ because $(1+\chi^6)-(\chi+\chi^5)=(1-\chi)(1-\chi^5)\ge0$.
So the quotient of the two terms of $\kappa_3$ is $\omega^{j-k}$ times a number of argument at
most $\pi/4$ in absolute value, hence is not $-1$.  Therefore $\kappa_3\ne0$, and each of these
labels is isolated by Lemma~\ref{lem:d2-isolation}(b).

\emph{Fold labels of rank $20$.}  At $1080$ fold labels the Jacobian of the twenty-two
polynomials has rank $20$.  Twenty rows together with $h=W_4-W_4(p)$ have an invertible
derivative: the rows of $F_{\mathrm{sel}}$ at $864$ labels, and all rows other than $F^Z_2$ and
$F^Z_3$ at the other $216$.  An omitted row has a nonzero coefficient of $h^3$ along the curve
these rows define, so each of these labels is isolated by Lemma~\ref{lem:d2-isolation}(c).

\emph{Fold labels of rank $19$.}  The other $72$ fold labels form one orbit under the maps
listed below.  Its seed $p_*=p_*(\nu,M)$ has the pair part of Lemma~\ref{lem:d2-fibre}(a) with
$D_{12}=M^{-1/3}$ and the fold triangle
\[
  (W_3,W_4,W_5)=(\omega,\omega^2,1),\qquad Z_B=-\delta_*W_B,\qquad
  (D_{34},D_{35},D_{45})=\delta_*(1,\omega,\omega^2).
\]
Apply Lemma~\ref{lem:d2-corank-two} with $F_{R'}=(F^W_2,F^W_4,F^W_5)$, $F_R$ the other nineteen
rows and $y=Z_3$.  The kernel coordinates are then the chart parameter $t$ and
$a=Z_3-Z_3(p_*)$.  At $t=0$ the derivative of $(F_R,t,Z_3)$ is block lower triangular, as in the
proof of Lemma~\ref{lem:d2-second}.  Its diagonal blocks are $(1)$ for $dt$; the triangle block
(rows $F^Z_3,F^W_3,F^Z_4,F^Z_5,T_{34},T_{35},T_{45},G_2,dZ_3$, columns $Z_B,W_B,D_{BB'}$), with
determinant $18\delta_*$; the pair block (rows $F^Z_1,F^W_1,F^Z_2,T_{12},G_1$, columns
$Z_A,W_A,D_{12}$), with determinant $-3M^{-1/3}$; and the diagonal crossed block $(2/D_{AB})$,
with determinant $-512i\,\nu^3(1-\nu^2)^3/(1+\nu^2)^6$.  So the invertibility hypothesis holds
for all $(\nu,M)$.  The differentials of the three omitted rows vanish on the two-dimensional
kernel of $dF_R(p_*)$, so (a) holds.  Write $g=(\psi_2,\psi_4,\psi_5)$ with
$\psi_k=F^W_k\circ\phi$.

At $M=1$ the jet of $g$ through degree four was computed exactly over
$\mathbb Q(\nu)(\sqrt3,i,\delta_*)$, and the nineteen rows of $F_R$ vanish through degree four on
the computed jet of $\phi$.  In the quadratic part, $\psi_2^{(2)}=0$, there is no $a^2$ term,
and the $3\times3$ coefficient matrix has rank one.  So $g^{(2)}=\psi_4^{(2)}\,n$ for a fixed
vector $n=(0,1,n_5)$, where $\psi_4^{(2)}=t(\gamma_{11}a+\gamma_{02}t)$ and
\[
  \gamma_{11}=\frac{9i\delta_*(\nu-i)(\nu+i)^3\,\Upsilon(\nu)}{4\nu(\nu^2-1)(\nu^2+1)^2},\qquad
  \Upsilon(\nu)=\nu^4-2i\nu^3-6\nu^2+2i\nu+1=\tfrac14\bigl(3(\nu-i)^4+(\nu+i)^4\bigr).
\]
For real $\nu$, $|3(\nu-i)^4|=3(\nu^2+1)^2>|(\nu+i)^4|$, so $\Upsilon(\nu)\ne0$ and
$\gamma_{11}\ne0$ for $0<\nu<1$.  Hence the common zeros of $g^{(2)}$ are two lines: the
\emph{boundary line} $\xi_\partial=(1,0)$, tangent to $\{t=0\}$, and the \emph{slanted line}
$\xi_s=(-\gamma_{02},\gamma_{11})$.  Moreover $\Pi_\xi=n\,d\psi_4^{(2)}(\xi)$, with
$d\psi_4^{(2)}(\xi_\partial)=(0,\gamma_{11})$ and
$d\psi_4^{(2)}(\xi_s)=\gamma_{11}(\gamma_{11},\gamma_{02})$, both nonzero.  So (b) holds.
On the slanted line take $N=3$: the matrix $(\Pi_{\xi_s}\,|\,r_3(\xi_s))$ has the $2\times2$
minor
\[
  \mu_3=-\frac{3^{10}(\nu-i)^4(\nu+i)^{16}\,\Upsilon(\nu)^5}
             {2^{12}\delta_*\,\nu^6(\nu^2-1)^6(\nu^2+1)^8}\ne0 .
\]
On the boundary line take $e=(0,1)$, so that $\varpi=\gamma_{11}n$, and $N=4$.  Identically in
$\nu$, $r_3(\xi_\partial)=g^{(3)}(1,0)=-b_2\varpi$, and the matrix
$(\Pi_{\xi_\partial}\,|\,r_4(\xi_\partial))$ has the minor
\[
  \mu_4=-\frac{27i\delta_*(\nu-i)^3(\nu+i)^5\,\Upsilon(\nu)}{64\,\nu^3(\nu^2-1)^3}\ne0 .
\]
So (c) holds.  Every coefficient used is regular for $0<\nu<1$: the jets are polynomials in the
coordinates of $p_*$, the masses and the entries of the inverses of the three diagonal blocks,
and the only further division is by $\gamma_{11}$.  Hence the identities above hold at every
$\nu\in(0,1)$.

For general $M$, $p_*(\nu,M)$ is obtained from $p_*(\nu,1)$ by multiplying $W_1,W_2$ by $M^{2/3}$
and $D_{12}$ by $M^{-1/3}$.  The map
$\tilde g(a,\tilde t)=\bigl(M^{-2/3}\psi_2,\psi_4,\psi_5\bigr)(a,\tilde t/M)$ was computed in
exact Laurent arithmetic in $M^{1/3}$, and the selected rows vanish along both lines through the
orders used.  The quadratic part of $\tilde g$ does not depend on $M$, the compatibility
$r_3(\xi_\partial)\in\C\varpi$ persists, and the two minors become $M^{-3}\mu_3(\nu)$ and
$M^{-5/3}\mu_4(\nu)$.  By Lemma~\ref{lem:d2-corank-two}, applied to $\tilde g$, $p_*(\nu,M)$ is
an isolated point of $\mathcal X_m$ for every $0<\nu<1$ and $M>0$.

\emph{Transport.}  Two involutions preserve the twenty-two polynomials, since these have real
coefficients and $m_3=m_4$: the reversal $(t,D_{AB})\mapsto(-t,-D_{AB})$, which is biholomorphic,
and complex conjugation followed by the exchange of bodies $3$ and $4$, which is
antiholomorphic.  They commute, fix $(\nu,M)$, and carry $p_*$ to the labels with the other three
signs of the orbit.  Multiplying $t$, $Z_B$, $W_A$, $D_{12}$ and $D_{BB'}$ by $\omega^j$
multiplies each polynomial by a cube root of unity and fixes the gauges.  A permutation of the
bodies $3,4,5$, followed by the element of the chart torus that restores the gauges, is a local
biholomorphism of $\mathcal X_m$ by the implicit function theorem, because the transversality
determinant is $2(Z_2-Z_1)W_5=2$.  These $3\cdot6$ maps give eighteen labels for each of the four
signs.  The resulting $72$ labels are exactly one orbit of the discrete label action and are
disjoint from the $1080$ labels above.  All these maps preserve isolatedness.

The computations in the proof of Lemma~\ref{lem:d2-L5}, namely the coefficients $\kappa_3$, the
ranks and the coefficients of $h^3$ at the $1080$ fold labels, the determinants, jets and minors
at $p_*$, and the orbit of $p_*$, are exact, over number fields containing $\omega$, $i$ and
$\delta_*$ and over rational function fields over them in $\nu$ and $M^{1/3}$
(Appendix~\ref{app:verification}).

\begin{lemma}[computer-assisted]\label{lem:d2-rays}
\textup{(a)} The zero set of $L_8$ in $\Rp^3$ is $\Rp(p_*,q_*,1)\cup\Rp(q_*,p_*,1)$ with
$p_*=0.67987758\ldots$ and $q_*=0.82343785\ldots$.
\textup{(b)} The zero set of $L_1$ in $\Rp^3$ is $\Rp(a_*,b_*,1)\cup\Rp(a_*,1,b_*)$ with
$a_*=1.02227400\ldots$ and $b_*=0.90559635\ldots$; at $(a_*,b_*,1)$ the zeros of $S_1$ that
satisfy \textup{(25)} form exactly one pair of complex conjugate non-real shapes.
\textup{(c)} Consequently the zero sets of $L_1L_2L_3$ and of $L_6L_7L_8$ in $\Rp^3$ are the rays
through the six permutations of $(a_*,b_*,1)$ and of $(1,p_*,q_*)$, respectively.
\end{lemma}
For (a) we combine exact Bernstein subdivision of the mass simplex with the Krawczyk method in
rational interval arithmetic.  For (b) we eliminate the shape by subresultants, locate the zeros
by Bernstein subdivision and the signs of resultants, enclose them by the Krawczyk method, and
show that the quadratic subresultant has negative discriminant there
(Appendix~\ref{app:verification}).  Part (c) follows from (a), (b) and the permutation statement
after Definition~\ref{def:d2-mass}.

\begin{lemma}[computer-assisted]\label{lem:d2-Lray}
Let $m\in\Rp^5$ with $m'$ on one of the rays of Lemma~\ref{lem:d2-rays}(c), and let $p$ be a
same-sign label at $m$ with $\Phi(p)=0$.  Then $J_T(p)$ is invertible and $\Psi_T(p)\ne0$, so $p$
is isolated at second order.
\end{lemma}

At each permutation of $(1,p_*,q_*)$ the $48$ Lagrange triangles, and at each permutation of
$(a_*,b_*,1)$ the six Euler triangles on the conjugate shape pair, are enumerated; there only the
ray factor itself vanishes among $L_1,\dots,L_8,L_5'$, so by Proposition~\ref{prop:d2-first}(a)
no other triangle carries a same-sign label with $\Phi=0$.  For $36+36$ triangles a sign vector
with $\sum_Bm_Bs_Bv_B=0$ exists, unique up to sign; $J_T$ is invertible, and complex interval
arithmetic gives $|\Psi_T|>0.36$ ($L_6$--$L_8$) and $|\Psi_T|>0.3$ ($L_1$--$L_3$).  Common mass
scaling covers the rays, and $\Psi_T$ does not involve $m_1,m_2$.
The shapes and the square and cube roots are enclosed by the Krawczyk method, and the solutions
of the linear systems are enclosed rigorously (Appendix~\ref{app:verification}).

\begin{lemma}[computer-assisted]\label{lem:d2-K}
Let $m\in\Rp^5$ with $m_1=m_2$, let $p$ be a one-different label at $m$ with odd body $5$ and
$\Phi(p)=0$, and let $S_j$ carry its triangle, so that $K_j(m')=0$.  Then $p$ is an isolated
point of $\mathcal X_m$.  More precisely:
\begin{alphalist}
  \item $j\notin\{3,6,7,8\}$, by Lemma~\ref{lem:d2-empty};
  \item if $j\in\{1,2\}$: at every zero of $K_1$ in $\Rp^3$ the $S_1$ shapes compatible with
  \textup{(28)} at $y=0$ reduce to one simple common root, which is real and lies on one of two
  explicit real branches; $J_T(p)$ is invertible and $q'(0)\ne0$.  The case $j=2$ is the image of
  $j=1$ under the exchange of bodies $3$ and $4$;
  \item if $j=4$: $m'$ lies on one of four arcs, $k\mapsto\mathrm{II}(k),\mathrm{III}(k)$ for
  $5^{-1/4}<k<1$ and their images under $m_3\leftrightarrow m_4$, where
  $\alpha=4k(1+k^4)$, $\beta=(1-k^2)(1+6k^2+k^4)$, $\gamma=(1+k^2)(6k^2-1-k^4)$,
  $\mathrm{II}=((\gamma^3-\beta^3)\alpha,(\alpha^3-\gamma^3)\beta,(\alpha^3-\beta^3)\gamma)$ and
  $\mathrm{III}=((\gamma^3-\beta^3)\alpha,(\gamma^3+\alpha^3)\beta,(\alpha^3+\beta^3)\gamma)$;
  $J_T(p)$ is invertible and $q'(0)\ne0$;
  \item if $j=5$: $m'\in\Rp(1,1,\rho_*)$, where $\rho_*\in(1.5736409,1.573641)$ is the unique
  positive root of $\rho^8+2\rho^7+3\rho^6-4\rho^5-5\rho^4-6\rho^3-3\rho^2-12\rho-12$, and every
  label at $m$ with a Lagrange triangle is isolated, for all $27$ classes and all $64$
  crossed-sign assignments.
\end{alphalist}
\end{lemma}

In (b) and (c), $60$ of the $64$ crossed-sign assignments have $\Phi\ne0$, and for the other four
(one orbit of conjugation and $t\mapsto-t$) $q'(0)\ne0$.  In (d), the two-equal classes have
$\Phi\ne0$; for the all-equal classes $60$ assignments have $\Phi\ne0$ and $4$ have $q'(0)\ne0$;
for the fold classes an augmented system in $h=W_4-W_4(p)$ gives a nonzero coefficient of order
three ($60$ assignments) or four ($4$).  Exact grading makes all items uniform in $m_1=m_2$.
In (b) the common root is obtained from a linear subresultant over $\mathbb Q$, using the
smoothness of the zero set of $K_1$ in $\Rp^3$, which is proved with Bernstein coefficients; the
two real branches are found from a rational parametrization of $S_1$ and Sturm sequences, and
$q'(0)\ne0$ follows from an exact reduction and an interval cover of the branch parameter.  In (c)
the computation is exact, on the rational parametrizations of the four arcs, and classifies the
real $S_4$ triangles; by Lemma~\ref{lem:d2-empty}(a) every $S_4$ triangle at positive masses is
real.  In (d) the zero set of $K_5$ in $\Rp^3$ is determined by exact resultants and a Sturm
count, and the labels are classified in complex rational interval arithmetic, at one triangle in
each orbit of conjugation and the exchange of bodies $3$ and $4$
(Appendix~\ref{app:verification}).  Renumbering the bodies as at the beginning of this section
gives the same conclusions for all pair splits and odd-body placements.

\begin{lemma}[computer-assisted]\label{lem:d2-P}
Let $m\in\Rp^5$ with $m_1\ne m_2$, and let $p$ be a one-different label at $m$ with odd body $5$,
$\Phi(p)=0$, whose triangle is not a fold triangle.  Then $J_T(p)$ is invertible and
$\Psi_1(p)\ne0$, so $p$ is isolated at second order.
\end{lemma}

We cover the simplex $m_3+m_4+m_5=1$ for $S_1,S_3,S_4$ (all Euler shapes, enclosed
by Krawczyk or Pellet discs) and $S_5,S_6,S_8$ (all Lagrange shapes, up to conjugation and the
cube-root action), for both relative signs $s_3s_4$, by boxes on each of which $G$ has a strict
sign, or $\Psi_1$ is enclosed away from $0$ in ball arithmetic, or, on boxes meeting $m_3=0$ or
$m_4=0$ (where $\Psi_1$ vanishes), its transversal derivative is.  The sets left out of the cover
are treated separately: the vertex corners ($G\ne0$); the ends $(m_3,m_4)=(0,\frac12)$ of $S_1$
and $S_4$ and $(\frac12,0)$ of $S_4$, where a majority body becomes massless (an exact edge
identity and interval margins $0.787$, $0.588$, $0.551$); the corner $(\frac12,\frac12,0)$ of
$S_3$ (an exact weighted blow-up of weight two in $m_5$ of the Euler shape parameter, after which
the determinant and the numerator of the response are $m_5^{13}$ and $m_5^6$ times factors bounded
below in modulus by $6.96\cdot10^8$ and $1.63\cdot10^6$ in interval arithmetic); six real points
of $S_4$ where the Euler label degenerates, and the corner $(\frac12,\frac12,0)$ of $S_4$ (an
exact closed form of $\Psi_1$ on the Euler shapes, in which the degenerating factor cancels); and
$S_5$ by exact algebra (Appendix~\ref{app:verification}).  On $S_5$ let $g$ be the
$W_4$-component of the response of the triangle block (row $F^Z_5$ omitted, $dW_5=0$) to a unit
forcing of the row $F^Z_4$, and $X=\tilde v_4/W_4-\tilde v_3/W_3$.  Then
$\Psi_1=-\frac12m_4gX^2$; the numerator of $g$ is a nonzero constant times
$m_3(m_3^2+m_3m_4+m_4^2)^2(P_0+\omega P_1)$ with real polynomials $P_0,P_1$ whose only common zero
in the simplex is $m_3=m_4=\frac13$, and $X\ne0$ except at the same point.  There $m_3=m_4=m_5$,
and the label with $G=0$ has $\rho=-1$, which is impossible.  $S_2$ and $S_7$ follow from $S_1$
and $S_6$ by exchanging bodies $3$ and $4$.

\begin{lemma}[computer-assisted]\label{lem:d2-fold}
Let $m\in\Rp^5$ and let $p$ be a one-different label at $m$ with odd body $5$, $\Phi(p)=0$ and a
fold triangle.  Then $dF_{\mathrm{sel}}(p)$ has rank $20$ and $Q|_C$ has a simple zero at $p$; in
particular $p$ is excluded.
\end{lemma}

$J_T(p)$ has corank exactly one, with an explicit kernel vector and left null vector $\ell$.
Since the $t$-column of the triangle rows is $-\varphi_B/W_B$ on the $Z$-rows, $\ell$ applied to it
is $c_\sigma X$ with $X$ as in Lemma~\ref{lem:d2-P}, and $X\ne0$ at every label considered (the
exceptional point $m_3=m_4=m_5$ carries no one-different label with $\Phi=0$ at positive masses).
So $C$ exists, is parametrized by $h=W_4-W_4(p)$, has $dt/dh=0$ at $p$, and
$dQ/dh=-\frac12c_\sigma m_4X\ne0$ at $p$.
The proof is an exact computation over $\mathbb Q(m_3,m_4)(\omega)$, with $m_3,m_4$ as
indeterminates: it gives the kernel vector, the left null vector, a minor of order eight of
$J_T(p)$ that does not vanish on the simplex, and the zero set of $X$
(Appendix~\ref{app:verification}).

\begin{theorem}\label{thm:d2}
No Diagram~2 germ exists at positive masses: for $m\in\Rp^5$ no curve in $A_m$ has a Diagram~2
germ.  In particular Theorem~\ref{thm:exclusion}\textup{(a)} holds for $r=2$.
\end{theorem}

\begin{proof}
Let $\gamma$ be a Diagram~2 germ on a curve in $A_m$ and $p$ its limit label
(Proposition~\ref{prop:d2-cover}).  For each $B$, $\sigma_Bx$ is the limit of
$m_2\delta_{2B}/(m_1\delta_{1B})$ along $\gamma$, which does not depend on the chart; so the
crossed pattern of $p$ is determined by $\gamma$, and renumbering the bodies $3,4,5$ before
applying Proposition~\ref{prop:d2-cover} permutes it.  We may therefore assume that $p$ is
same-sign or has odd body $5$.
By Lemma~\ref{lem:d2-isolation}(a) it suffices to show that $p$ is excluded.  By
Proposition~\ref{prop:d2-first}, $\Phi(p)=0$; let $S_j$ carry the triangle of $p$
(Lemma~\ref{lem:d2-fibre}).

\emph{Same-sign.}  By Proposition~\ref{prop:d2-first}(a), $L_j(m')=0$, or $L_5L_5'=0$ if
$j=5$.  Lemma~\ref{lem:d2-empty} excludes $j=4$ and $L_5'=0$.  If $L_5=0$, then $m_3=m_4=m_5$,
$j=5$ by Lemma~\ref{lem:d2-empty}(b), and Lemma~\ref{lem:d2-L5} applies.  Otherwise
$j\in\{1,2,3,6,7,8\}$, $m'$ lies on a ray of Lemma~\ref{lem:d2-rays}(c), and
Lemma~\ref{lem:d2-Lray} applies.

\emph{One-different.}  If $m_1=m_2$, Lemma~\ref{lem:d2-K} applies.  If $m_1\ne m_2$,
Lemma~\ref{lem:d2-fold} applies when the triangle is a fold triangle and Lemma~\ref{lem:d2-P}
otherwise.

In every case $p$ is excluded (by Definition~\ref{def:d2-psi} for the labels isolated at second
order), which contradicts Lemma~\ref{lem:d2-isolation}(a).
\end{proof}

\begin{remark}\label{rem:d2-scope}
(i) Theorem~\ref{thm:d2} is stronger than Theorem~\ref{thm:exclusion}(a) requires: it excludes
every Diagram~2 germ, with or without an unbounded four-product, and uses no reality.  From
\cite{AK} it uses only Fact~\ref{fact:rules} (in Proposition~\ref{prop:d2-cover}), besides
Lemma~\ref{lem:germ}; the mass conditions (25)--(28) of \cite[\S8.2]{AK} are derived
(Proposition~\ref{prop:d2-first}), not assumed.

(ii) One can show that the layer $u=\infty$ of the $P_j$ misses the positive cone, and that no
non-degenerate triple of positions with positive centre-of-mass masses lies on two curves
$S_i\ne S_j$ (Appendix~\ref{app:verification}).  We do not use this.
\end{remark}
